% ============================================================
%  Impacto da IA em Repositórios GitHub — Lab04 / UFMG
%  Compilar no Overleaf (pdflatex) — pacotes padrão CTAN
% ============================================================
\documentclass[12pt,a4paper,twocolumn]{article}

\usepackage[utf8]{inputenc}
\usepackage[T1]{fontenc}
\usepackage[brazil]{babel}
\usepackage{geometry}
\usepackage{graphicx}
\usepackage{booktabs}
\usepackage{tabularx}
\usepackage{multirow}
\usepackage{pgfplots}
\usepackage{pgfplotstable}
\usepackage{tikz}
\usepackage{xcolor}
\usepackage{hyperref}
\usepackage{caption}
\usepackage{subcaption}
\usepackage{amsmath}
\usepackage{setspace}
\usepackage{microtype}
\usepackage{enumitem}
\usepackage{float}
\usepackage{array}
\usepackage{parskip}
\usepackage{titlesec}
\usepackage{fancyhdr}
\usepackage{abstract}
\usepackage{colortbl}
\usepackage{authblk}

\pgfplotsset{compat=1.18}

% ── Geometria ──────────────────────────────────────────────
\geometry{a4paper, top=2.4cm, bottom=2.4cm, left=1.8cm, right=1.8cm,
          columnsep=0.7cm}

% ── Cores ──────────────────────────────────────────────────
\definecolor{primaryblue}{RGB}{30,100,200}
\definecolor{accentcyan}{RGB}{0,180,210}
\definecolor{accentgreen}{RGB}{0,170,120}
\definecolor{accentred}{RGB}{210,55,70}
\definecolor{accentamber}{RGB}{200,140,10}
\definecolor{lightgray}{RGB}{245,247,250}
\definecolor{darkgray}{RGB}{80,100,120}

% ── Hyperref ───────────────────────────────────────────────
\hypersetup{colorlinks=true, linkcolor=primaryblue,
            citecolor=primaryblue, urlcolor=accentcyan,
            pdftitle={Impacto da IA em Repositorios GitHub}}

% ── Seções ─────────────────────────────────────────────────
\titleformat{\section}{\large\bfseries\color{primaryblue}}
            {\thesection.}{0.5em}{}[\vspace{-4pt}\rule{\columnwidth}{0.4pt}]
\titleformat{\subsection}{\normalsize\bfseries}{\thesubsection.}{0.5em}{}
\titlespacing{\section}{0pt}{10pt}{4pt}
\titlespacing{\subsection}{0pt}{6pt}{2pt}

% ── Cabeçalho/Rodapé ───────────────────────────────────────
\pagestyle{fancy}
\fancyhf{}
\fancyhead[L]{\small\color{darkgray}\textit{Lab04 --- Impacto da IA em Repositórios GitHub}}
\fancyhead[R]{\small\color{darkgray}2025}
\fancyfoot[C]{\small\thepage}
\renewcommand{\headrulewidth}{0.4pt}
\renewcommand{\headrule}{\hbox to\headwidth{\color{accentcyan}\leaders\hrule height\headrulewidth\hfill}}

% ── Pgfplots style global ──────────────────────────────────
\pgfplotsset{
  dashstyle/.style={
    width=\columnwidth, height=5.5cm,
    ybar, bar width=10pt,
    enlarge x limits=0.3,
    symbolic x coords={Mediana, Media},
    xtick=data,
    ymin=0,
    grid=major, grid style={dotted, gray!40},
    legend style={font=\tiny, at={(0.5,1.02)}, anchor=south,
                  legend columns=-1, draw=none, fill=none},
    tick label style={font=\tiny},
    label style={font=\tiny},
    title style={font=\small\bfseries, at={(0.5,1.05)}, anchor=south},
  },
  groupstyle/.style={
    width=\columnwidth, height=5.2cm,
    ybar, bar width=12pt,
    enlarge x limits=0.25,
    symbolic x coords={Baixo, Medio, Alto},
    xtick=data,
    grid=major, grid style={dotted, gray!40},
    legend style={font=\tiny, at={(0.5,1.02)}, anchor=south,
                  legend columns=-1, draw=none, fill=none},
    tick label style={font=\tiny},
    label style={font=\tiny},
    title style={font=\small\bfseries, at={(0.5,1.05)}, anchor=south},
  }
}

% ═══════════════════════════════════════════════════════════
\begin{document}
% ═══════════════════════════════════════════════════════════

% ── Título (onecolumn temporariamente) ────────────────────
\twocolumn[{%
\begin{@twocolumnfalse}
\vspace*{0.4cm}
{\centering
  {\LARGE\bfseries Impacto de Ferramentas de Inteligência Artificial\\[4pt]
   em Repositórios GitHub: Uma Análise Empírica\par}
  \vspace{0.5em}
  {\large Andre Lazarini\par}
  \vspace{0.2em}
  {\small Laboratório de Experimentação em Engenharia de Software\par}
  {\small \href{mailto:andrexlazarini@gmail.com}{andrexlazarini@gmail.com}\par}
  \vspace{0.3em}
  {\small Junho de 2025\par}
\par}

\vspace{0.6em}
\noindent\rule{\linewidth}{0.5pt}

% ── Resumo (2 idiomas) ────────────────────────────────────
\begin{center}\textbf{Resumo}\end{center}
\small
Ferramentas de inteligência artificial (IA) tornaram-se parte do fluxo de trabalho
de muitos projetos de software de código aberto. Este trabalho apresenta uma análise
empírica de 105 repositórios Python públicos do GitHub, classificados por nível de
evidência de uso de IA. Comparamos métricas de atividade de manutenção e de
manutenibilidade nos períodos anterior e posterior ao primeiro sinal identificado de
uso de IA. Os resultados indicam redução nas medianas de commits, \textit{issues} e
tempo de resolução após a adoção de IA, e um resultado misto para o
Maintainability Index (MI), sem evidência de melhora uniforme de manutenibilidade.

\vspace{0.4em}
\noindent\textbf{Palavras-chave:} mineração de repositórios, inteligência artificial,
engenharia de software empírica, manutenibilidade.

\vspace{0.5em}
\begin{center}\textbf{Abstract}\end{center}
Artificial intelligence (AI) tools have become part of the workflow of many open-source
software projects. This paper presents an empirical analysis of 105 public Python
repositories from GitHub, classified by level of AI usage evidence. We compare
maintenance activity and maintainability metrics in the periods before and after the
first identified AI usage signal. Results indicate reductions in the medians of commits,
issues, and resolution time after AI adoption, and a mixed result for the Maintainability
Index (MI), with no clear uniform maintainability improvement.

\noindent\textbf{Keywords:} repository mining, artificial intelligence,
empirical software engineering, maintainability.

\noindent\rule{\linewidth}{0.5pt}
\vspace{0.6em}
\end{@twocolumnfalse}
}]

% ═══════════════════════════════════════════════════════════
\section{Introdução}
% ═══════════════════════════════════════════════════════════

Ferramentas de IA generativa --- como GitHub Copilot, Claude e ChatGPT --- passaram a
apoiar tarefas de codificação, revisão e documentação em projetos de software. Com isso,
surge a questão de como essas ferramentas afetam o ciclo de vida do software, em especial
métricas de atividade de manutenção e de qualidade do código.

Estudos recentes em mineração de repositórios (MSR) investigaram o uso de
\textit{co-authorship} por IA em \textit{commits}~\cite{dakhel2023copilot},
bem como o impacto de ferramentas como Copilot na produtividade de
desenvolvedores~\cite{ziegler2022copilot}. Entretanto, há lacuna em análises
empíricas de larga escala que comparem métricas \textit{antes} e \textit{depois}
da introdução de sinais de IA em repositórios públicos.

Este trabalho responde às seguintes questões de pesquisa:

\begin{itemize}[leftmargin=1.2em, itemsep=2pt]
  \item \textbf{RQ1:} Como as métricas de atividade de manutenção mudam após a
    identificação de sinais de uso de IA nos repositórios?
  \item \textbf{RQ2:} Como as métricas de qualidade e manutenibilidade mudam após a
    identificação de sinais de uso de IA nos repositórios?
\end{itemize}

% ═══════════════════════════════════════════════════════════
\section{Trabalhos Relacionados}
% ═══════════════════════════════════════════════════════════

Dakhel et al.~\cite{dakhel2023copilot} avaliaram a qualidade do código gerado pelo
GitHub Copilot e identificaram limitações em lógica complexa. Ziegler et al.
\cite{ziegler2022copilot} mediram a taxa de aceitação de sugestões do Copilot por
desenvolvedores profissionais. Nguyen e Nadi~\cite{nguyen2022copilot} avaliaram
a corretude funcional de código gerado por IA.

Diferentemente desses trabalhos, nossa abordagem é \textit{observacional} e baseada em
mineração de repositórios reais, sem experimentos controlados, focando em mudanças
temporais nas métricas dos projetos após a introdução de sinais de IA.

% ═══════════════════════════════════════════════════════════
\section{Metodologia}
% ═══════════════════════════════════════════════════════════

\subsection{Coleta e Seleção de Repositórios}

A amostra foi constituída por 105 repositórios públicos Python do GitHub, selecionados
por relevância (número de estrelas $\geq$ 1.000) e presença de pelo menos um sinal de
uso de IA. Os dados foram coletados via API do GitHub e ferramentas de análise
estática de código (para o Maintainability Index).

\subsection{Identificação de Sinais de IA}

Quatro tipos de evidência foram investigados:

\begin{enumerate}[leftmargin=1.4em, itemsep=1pt]
  \item \textbf{Co-autoria em commits:} mensagens com \texttt{Co-Authored-By}
    referenciando modelos de IA;
  \item \textbf{Menção ao Claude:} citações ao modelo Anthropic Claude;
  \item \textbf{Menção ao Copilot:} citações ao GitHub Copilot;
  \item \textbf{Arquivo relacionado a IA:} \texttt{.cursorrules},
    \texttt{.github/copilot-instructions.md} e similares.
\end{enumerate}

Cada repositório recebe um \textit{AI Score} correspondente ao número total de
evidências encontradas. Os projetos foram agrupados em três níveis: \textbf{Baixo}
(score 2--3), \textbf{Médio} (4--6) e \textbf{Alto} (7+).

\subsection{Métricas Coletadas}

\textbf{RQ1 --- Atividade de manutenção:}
número de \textit{commits}, número de \textit{issues}, taxa de \textit{issues} por mês
e tempo médio de resolução de \textit{issues} (em horas).

\textbf{RQ2 --- Qualidade e manutenibilidade:}
percentual de \textit{commits} de correção (\textit{fixes}), percentual de
\textit{reverts}, linhas alteradas por \textit{commit} e
Maintainability Index (MI)~\cite{oman1992maintainability} médio e mediano.

\subsection{Método de Análise}

Para cada repositório, identificamos o primeiro sinal de IA (marco temporal) e
calculamos as métricas no período \textbf{anterior} e \textbf{posterior} a esse marco.
A mediana foi adotada como medida de tendência central para reduzir o efeito de
\textit{outliers} (repositórios muito grandes ou muito ativos).

% ═══════════════════════════════════════════════════════════
\section{Resultados}
% ═══════════════════════════════════════════════════════════

\subsection{Caracterização do Dataset}

A Tabela~\ref{tab:dataset} resume as características principais da amostra.

\begin{table}[H]
\centering
\caption{Caracterização geral da amostra.}
\label{tab:dataset}
\small
\begin{tabular}{lc}
\toprule
\textbf{Atributo} & \textbf{Valor} \\
\midrule
Repositórios         & 105 \\
Mediana de estrelas  & 17.169 \\
Mediana de idade     & 9,3 anos \\
Mediana AI Score     & 3 \\
Idioma principal     & Python \\
\bottomrule
\end{tabular}
\end{table}

A Figura~\ref{fig:grupos} apresenta a distribuição dos repositórios nos três
grupos de adoção de IA.

\begin{figure}[H]
\centering
\begin{tikzpicture}
\begin{axis}[
  width=0.9\columnwidth, height=4.2cm,
  ybar,
  bar width=22pt,
  symbolic x coords={Baixo (2-3), Medio (4-6), Alto (7+)},
  xtick=data,
  xticklabel style={font=\footnotesize},
  ymin=0, ymax=65,
  ytick={0,10,20,30,40,50,60},
  ylabel={\footnotesize Repositórios},
  ylabel style={font=\footnotesize},
  nodes near coords,
  nodes near coords style={font=\footnotesize},
  grid=major, grid style={dotted, gray!40},
  title={\footnotesize Repositórios por grupo de adoção de IA},
  title style={font=\small\bfseries, at={(0.5,1.05)}, anchor=south},
  enlarge x limits=0.2,
]
\addplot[fill=accentamber!80, draw=none] coordinates {
  (Baixo (2-3), 55) (Medio (4-6), 32) (Alto (7+), 18)
};
\end{axis}
\end{tikzpicture}
\caption{Distribuição dos 105 repositórios por grupo de adoção de IA.}
\label{fig:grupos}
\end{figure}

A Tabela~\ref{tab:grupos} detalha cada grupo com estatísticas de popularidade
e características de IA.

\begin{table}[H]
\centering
\caption{Estatísticas por grupo de adoção de IA.}
\label{tab:grupos}
\small
\setlength{\tabcolsep}{4pt}
\begin{tabular}{lrrrr}
\toprule
\textbf{Grupo} & \textbf{Repos} & \textbf{Med. $\bigstar$} & \textbf{Med. Idade} & \textbf{Med. Score} \\
\midrule
Baixo (2--3) & 55 & 14.662 & 9,50 a & 3,0 \\
Médio (4--6) & 32 & 18.905 & 9,00 a & 6,0 \\
Alto (7+)    & 18 & 30.102 & 9,05 a & 8,5 \\
\bottomrule
\end{tabular}
\end{table}

Repositórios com maior adoção de IA tendem a ser mais populares (mediana de estrelas
cresce de 14.662 no grupo Baixo para 30.102 no grupo Alto), o que pode refletir maior
visibilidade e adoção precoce de ferramentas modernas.

\subsection{RQ1 --- Atividade de Manutenção}

A Tabela~\ref{tab:rq1} e a Figura~\ref{fig:rq1} apresentam as comparações das
medianas de atividade antes e depois dos sinais de IA.

\begin{table}[H]
\centering
\caption{Métricas de atividade — medianas antes e depois.}
\label{tab:rq1}
\small
\setlength{\tabcolsep}{4pt}
\begin{tabular}{lrrr}
\toprule
\textbf{Métrica} & \textbf{Antes} & \textbf{Depois} & \textbf{$\Delta$} \\
\midrule
Commits                   & 255,00   & 166,00   & \textcolor{accentred}{$-$89,00} \\
Issues                    & 399,00   & 113,00   & \textcolor{accentred}{$-$286,00} \\
Taxa issues/mês           &  66,50   &  18,83   & \textcolor{accentred}{$-$47,67} \\
Tempo resolução (h)       & 9.898,01 & 5.645,14 & \textcolor{accentred}{$-$4.252,87} \\
\bottomrule
\end{tabular}
\end{table}

\begin{figure}[H]
\centering
\begin{tikzpicture}
\begin{axis}[
  dashstyle,
  height=5.2cm,
  ylabel={\footnotesize Commits (mediana)},
  title={\footnotesize RQ1 --- Commits: antes vs.\ depois},
  ymin=0, ymax=450,
  legend entries={Antes, Depois},
]
\addplot[fill=accentcyan!70, draw=none] coordinates {
  (Mediana, 255) (Media, 401)
};
\addplot[fill=accentgreen!70, draw=none] coordinates {
  (Mediana, 166) (Media, 365)
};
\end{axis}
\end{tikzpicture}
\begin{tikzpicture}
\begin{axis}[
  dashstyle,
  height=5.2cm,
  ylabel={\footnotesize Issues (mediana)},
  title={\footnotesize RQ1 --- Issues: antes vs.\ depois},
  ymin=0, ymax=1600,
  legend entries={Antes, Depois},
]
\addplot[fill=accentcyan!70, draw=none] coordinates {
  (Mediana, 399) (Media, 1406)
};
\addplot[fill=accentgreen!70, draw=none] coordinates {
  (Mediana, 113) (Media, 452)
};
\end{axis}
\end{tikzpicture}
\caption{Comparação de commits e issues antes e depois dos sinais de IA.}
\label{fig:rq1}
\end{figure}

As medianas de todas as métricas de atividade diminuem no período posterior.
Commits caem de 255 para 166 ($-$34,9\%), \textit{issues} de 399 para 113
($-$71,7\%) e tempo de resolução de 9.898~h para 5.645~h ($-$42,9\%).

\textbf{Resposta à RQ1:} as medianas indicam redução consistente no volume de atividade
de manutenção após a introdução de sinais de IA, incluindo menor tempo de resolução de
\textit{issues}.

\subsection{RQ2 --- Qualidade e Manutenibilidade}

A Tabela~\ref{tab:rq2} e a Figura~\ref{fig:rq2} apresentam as comparações
para as métricas de qualidade.

\begin{table}[H]
\centering
\caption{Métricas de qualidade — medianas antes e depois.}
\label{tab:rq2}
\small
\setlength{\tabcolsep}{4pt}
\begin{tabular}{lrrr}
\toprule
\textbf{Métrica} & \textbf{Antes} & \textbf{Depois} & \textbf{$\Delta$} \\
\midrule
\% Fixes            & 32,12 & 36,07 & \textcolor{accentgreen}{$+$3,95} \\
\% Reverts          &  1,10 &  1,24 & \textcolor{accentred}{$+$0,14} \\
Linhas/commit       & 193,79& 217,49& \textcolor{accentgreen}{$+$23,69} \\
MI médio            & 63,43 & 60,90 & \textcolor{accentred}{$-$2,53} \\
MI mediano          & 60,29 & 58,50 & \textcolor{accentred}{$-$1,80} \\
\bottomrule
\end{tabular}
\end{table}

\begin{figure}[H]
\centering
\begin{tikzpicture}
\begin{axis}[
  dashstyle,
  height=5.2cm,
  ylabel={\footnotesize \% Fixes (mediana)},
  title={\footnotesize RQ2 --- \% Fixes: antes vs.\ depois},
  ymin=0, ymax=45,
  legend entries={Antes, Depois},
]
\addplot[fill=accentcyan!70, draw=none] coordinates {
  (Mediana, 32.12) (Media, 33.45)
};
\addplot[fill=accentgreen!70, draw=none] coordinates {
  (Mediana, 36.07) (Media, 38.76)
};
\end{axis}
\end{tikzpicture}
\begin{tikzpicture}
\begin{axis}[
  dashstyle,
  height=5.2cm,
  ylabel={\footnotesize MI (mediana)},
  title={\footnotesize RQ2 --- Maintainability Index: antes vs.\ depois},
  ymin=55, ymax=70,
  legend entries={Antes, Depois},
]
\addplot[fill=accentcyan!70, draw=none] coordinates {
  (Mediana, 60.29) (Media, 63.43)
};
\addplot[fill=accentgreen!70, draw=none] coordinates {
  (Mediana, 58.50) (Media, 60.90)
};
\end{axis}
\end{tikzpicture}
\caption{Comparação do \% Fixes e do Maintainability Index antes e depois dos sinais de IA.}
\label{fig:rq2}
\end{figure}

O percentual mediano de \textit{fixes} aumenta de 32,12\% para 36,07\%, sugerindo
maior proporção de \textit{commits} de correção. Linhas alteradas por \textit{commit}
sobem de 193 para 217, indicando \textit{commits} mais expressivos. Porém, o MI
médio cai de 63,43 para 60,90 ($-$2,53), e o mediano de 60,29 para 58,50 ($-$1,80).

\textbf{Resposta à RQ2:} resultado misto. Há melhora em \textit{fixes} e tamanho de
\textit{commits}, mas o Maintainability Index apresenta pequena queda, sem evidência
de melhora uniforme de manutenibilidade.

% ═══════════════════════════════════════════════════════════
\section{Discussão}
% ═══════════════════════════════════════════════════════════

\subsection{Interpretação dos Resultados}

A queda nas métricas de atividade (RQ1) pode ter múltiplas explicações não excludentes:
(\textit{i}) a IA pode ter aumentado a eficiência, reduzindo a necessidade de abrir
\textit{issues} e acelerando a resolução; (\textit{ii}) o período posterior pode ser
naturalmente menor em termos de janela temporal disponível para acúmulo de eventos;
(\textit{iii}) projetos mais maduros tendem a ter atividade decrescente independentemente
da IA.

Para a RQ2, o aumento de \textit{fixes} e o crescimento de linhas por \textit{commit}
sugerem mudanças no estilo de desenvolvimento --- \textit{commits} maiores e com maior
proporção de correções --- possivelmente induzidos pela colaboração com IA. A queda no
MI, embora pequena, indica que o código gerado ou refatorado com IA pode ser
ligeiramente mais complexo ou menos legível pelos critérios da métrica.

\subsection{Análise por Grupos}

Repositórios do grupo Alto (score $\geq 7$) têm mediana de estrelas de 30.102, quase
o dobro do grupo Baixo (14.662). Isso pode indicar que projetos mais populares adotam
IA com mais intensidade, ou que ferramentas de IA são mais visíveis em projetos de
grande escala.

% ═══════════════════════════════════════════════════════════
\section{Ameaças à Validade}
% ═══════════════════════════════════════════════════════════

\textbf{Validade interna:} o desenho é observacional, sem grupo de controle.
Causalidade não pode ser inferida. Variáveis de confusão (mudança de equipe,
maturidade do projeto, etc.) não são controladas.

\textbf{Validade de construto:} o AI Score é um indicador de evidência, não de
intensidade de uso. Um repositório com alto score pode usar IA de forma superficial.
O MI é sensível à linguagem e às ferramentas de análise usadas.

\textbf{Validade externa:} a amostra é restrita a Python e a projetos com mais de
1.000 estrelas. Projetos menores ou em outras linguagens podem apresentar padrões
diferentes.

\textbf{Validade de conclusão:} a escolha da mediana mitiga outliers, mas não elimina
o efeito de projetos com comportamento atípico.

% ═══════════════════════════════════════════════════════════
\section{Conclusão}
% ═══════════════════════════════════════════════════════════

Este trabalho apresentou uma análise empírica do impacto de ferramentas de IA em 105
repositórios Python públicos do GitHub. Os principais achados são:

\begin{itemize}[leftmargin=1.2em, itemsep=2pt]
  \item \textbf{RQ1:} redução consistente nas medianas de commits ($-$34,9\%),
    \textit{issues} ($-$71,7\%), taxa mensal de \textit{issues} ($-$71,7\%) e tempo
    de resolução ($-$42,9\%) após sinais de IA.
  \item \textbf{RQ2:} resultado misto --- \% de \textit{fixes} e linhas por
    \textit{commit} aumentam, mas o Maintainability Index apresenta pequena queda.
  \item Repositórios com maior adoção de IA tendem a ser mais populares (maior
    mediana de estrelas).
\end{itemize}

Trabalhos futuros devem considerar grupos de controle, análise qualitativa de
\textit{commits} assistidos por IA e replicação em outras linguagens de programação.

% ── Referências ───────────────────────────────────────────
\bibliographystyle{plain}
\begin{thebibliography}{9}

\bibitem{dakhel2023copilot}
  A.~M. Dakhel, V.~Majdinasab, A.~Nikanjam, F.~Khomh, M.~C. Desmarais e Z.~Jiang,
  ``GitHub Copilot AI pair programmer: Asset or Liability?''
  \textit{Journal of Systems and Software}, 203:111806, 2023.

\bibitem{ziegler2022copilot}
  A.~Ziegler, E.~Kalliamvakou, X.~A. Li, A.~Rice, D.~Rifkin,
  S.~Simister, G.~Sittampalam e E.~Aftandilian,
  ``Productivity Assessment of Neural Code Completion.''
  In: \textit{Proc. MAPS}, 2022.

\bibitem{nguyen2022copilot}
  N.~Nguyen e S.~Nadi,
  ``An empirical evaluation of GitHub Copilot's code suggestions.''
  In: \textit{Proc. MSR}, pp.\ 1--5, 2022.

\bibitem{oman1992maintainability}
  P.~Oman e J.~Hagemeister,
  ``Metrics for assessing a software system's maintainability.''
  In: \textit{Proc. ICSM}, pp.\ 337--344, 1992.

\bibitem{kalliamvakou2014github}
  E.~Kalliamvakou et al.,
  ``The promises and perils of mining GitHub.''
  In: \textit{Proc. MSR}, pp.\ 92--101, 2014.

\end{thebibliography}

\end{document}
